\scanpage{89}
$2$ forment les sections d'un torseur sous $(\unsure{\mathfrak S|\mathcal G}_3)_S$. Donc les bi-épinglages forment les sections d'un torseur sous le groupe \deleted{\unsure{absolu|constant}}

\[
(\mathord{\mathrm{I}\mkern-7mu\mu}_4\times \mathfrak S_3)_S
\tag*{(45)}
\]

donc on trouve une action \inserted{\unsure{à g.|à d.}} du groupe absolu $(\mathord{\mathrm{I}\mkern-7mu\mu}_4\times \mathfrak S_3)_{S_0}$ sur \deleted{\unsure{U}} le schéma modulaire $(37)$ $U_{0,3}\times_{S_0}\mathord{\mathrm{I}\mkern-7mu\mu}_4^*$, qu'il faudrait expliciter, en $=$ temps que l'action sur

\[
\mathbb T_0=\underline O_{U_{0,3}\times \mathord{\mathrm{I}\mkern-7mu\mu}_4^*}(1)
\tag*{(46)}
\]

à l'aide de l'action ``tautologique'' de $\mathfrak S_3$ sur $\deleted{\mathbb P}U_{03}\unsure{=|\cdot}\Sigma_{\{0,1,\infty\}}\setminus\{0,1,\infty\}_{S_0}$ via l'action de $\mathfrak S_3$ sur $V_{\deleted{\unsure{\underline J|J}}\{0,1,\infty\}}(\underline O_{S_0})$.

\[
\left\{\begin{aligned}
&\mathfrak S_3\ \text{opère sur } U_{03}\times \mathord{\mathrm{I}\mkern-7mu\mu}_4^*\ \text{via }
\begin{cases}
\text{action tautol. sur }U_{03}\\
\text{le caractère }\operatorname{sg}(g)\text{ sur }\mathord{\mathrm{I}\mkern-7mu\mu}_4^*
\end{cases}\\[2pt]
&\mathord{\mathrm{I}\mkern-7mu\mu}_4\ \text{opère sur }U_{03}\times \mathord{\mathrm{I}\mkern-7mu\mu}_4^*\ \text{via }
\begin{cases}
\text{action triviale sur }U_{03}\\
\begin{gathered}\text{le car. can. }\mathord{\mathrm{I}\mkern-7mu\mu}_4\to \mathord{\mathrm{I}\mkern-7mu\mu}_2=\unsure{(\pm1)|\{\pm1\}}_{S_0}\\
\text{sur }\mu_4^*\end{gathered}
\end{cases}
\end{aligned}\right.
\tag*{(47)}
\]

donc $\mathfrak S_3\times \mathord{\mathrm{I}\mkern-7mu\mu}_4$ opère sur le facteur $U_{03}$ via l'

% archive sheet number: 88
